\section*{第 \ref{ch:b2:oceanic-islands} 章——洋岛}

\begin{solution}{exo:b2:oceanic-islands:1}
(a) 前\omterm{def:b2:oceanic-islands:stages}{盾状阶段}，如今的洛伊希\omterm{def:b2:oceanic-islands:seamount}{海山}即属此阶段。(b) \omterm{def:b2:oceanic-islands:stages}{盾状阶段}。(c) 后\omterm{def:b2:oceanic-islands:stages}{盾状阶段}。(d) 复活
阶段。(e) 不再有喷发阶段：火山已沉没于其礁体之下，该礁现为\omterm{def:b2:oceanic-islands:reefs}{环礁}。
\end{solution}

\begin{solution}{exo:b2:oceanic-islands:2}
(a) $2500 + 350\sqrt9 = \qty{3550}{m}$。(b) $350\times(\sqrt{25} - \sqrt9) = 350\times2 =
\qty{700}{m}$。(c) $350/(2\times5) = \qty{35}{m}$ 每百万年。
\end{solution}

\begin{solution}{exo:b2:oceanic-islands:3}
(a) $5000/17 = \qty{294}{ka}$，用于主要岛屿周围；$30\,000/70 = \qty{429}{ka}$，沿
海岭。(b) 滑坡未被计入：较年轻的火山可能掩埋较老火山的滑坡疤痕与碎屑，
较老的沉积层亦被沉积物覆盖。真实数目更大，真实间隔
更短。
\end{solution}

\begin{solution}{exo:b2:oceanic-islands:4}
(a) $4158\times0.3048 = \qty{1267}{m}$；$4610\times0.3048 = \qty{1405}{m}$；$150\times0.3048 =
\qty{46}{m}$。(b) 最初礁体的生长深度不超过 \qty{46}{m}，其底部现位于
\qty{1267}{m}，火山至少下沉 $1267 - 46 = \qty{1221}{m}$（若当时海平面更高，
则下沉量更大）。(c) 下沉量至多约 \qty{1267}{m}，历时至少 \qty{33.9}{Ma}：
即至多 $1267/33.9 = \qty{37}{m}$ 每百万年，即 \qty{0.037}{mm} 每年。
\end{solution}

\begin{solution}{exo:b2:oceanic-islands:5}
(a) $\dot S = 175/\sqrt t$：\qty{87.5}{m.Ma^{-1}}（\qty{0.0875}{mm.\text{年}^{-1}}）在 \qty{4}{Ma} 时，
\qty{35}{m.Ma^{-1}}（\qty{0.035}{mm.\text{年}^{-1}}）在 \qty{25}{Ma} 时，\qty{25}{m.Ma^{-1}}
（\qty{0.025}{mm.\text{年}^{-1}}）在 \qty{49}{Ma} 时。(b) $4.2/0.0875 = 48$ 倍。(c) $2.6/0.0875 = 30$
倍：加载阶段使岛屿下沉的速度远快于年轻板块冷却时的下沉。
\end{solution}

\begin{solution}{exo:b2:oceanic-islands:6}
(a) $\dot S = (3200/62.8)\,e^{-100/62.8} = 50.96\times0.2035 = \qty{10.37}{m}$ 每百万年，
\qty{0.0104}{m} 每千年，与实测的每千年 \qty{0.01}{m} 相符。(b) $e^{-t/62.8}
= 10/50.96$，$t = 62.8\ln5.096 = \qty{102.3}{Ma}$。(c) 平方根律：$2500 + 350\times5 =
\qty{4250}{m}$；指数律：$6400 - 3200\times0.6716 = \qty{4251}{m}$；二者相差
\qty{1}{m}。
\end{solution}

\begin{solution}{exo:b2:oceanic-islands:7}
(a) $\qty{8.6}{cm.\text{年}^{-1}} = \qty{86}{km.Ma^{-1}}$；$86\times33 = \qty{2838}{km}$。(b) 是的：
\qty{27.7}{Ma} 小于礁体淹没的时限 \qty{33}{Ma}，故该礁体
应仍存活。(c) $55\times1.609 = \qty{88.5}{km}$，再加 $88.5/86 = \qty{1.0}{Ma}$，约
\qty{28.7}{Ma}。此数略小于 \qty{33}{Ma}，但该数本身只是
对整条链的近似值，生长还取决于各\omterm{def:b2:oceanic-islands:reefs}{环礁}处
的水温。库雷\omterm{def:b2:oceanic-islands:reefs}{环礁}是最后一座仍存活的\omterm{def:b2:oceanic-islands:reefs}{环礁}，确实接近\omterm{def:b2:oceanic-islands:darwin-point}{达尔文点}。
\end{solution}

\begin{solution}{exo:b2:oceanic-islands:8}
(a) 长方体：$30\times17\times2 = \qty{1020}{km^3}$；棱锥为其三分之一，\qty{340}{km^3}。(b) 单独一块
仅为基拉韦厄火山整体的十分之一（长方体）或三十分之一（棱锥）。(c) $23\,000\times
0.1 = \qty{2300}{km^3}$，为基拉韦厄火山的四分之一；此估计成立与否取决于所假设的
厚度。
\end{solution}

\begin{solution}{exo:b2:oceanic-islands:9}
(a) $1000/2.6 = \qty{385}{ka}$。(b) 是的：\qty{2.6}{m} 每千年即为
\qty{2.6}{mm.\text{年}^{-1}}，小于 \qty{4.2}{mm.\text{年}^{-1}}，余量为 \qty{1.6}{mm.\text{年}^{-1}}。
(c) 相对上升速率 $12 + 2.6 = \qty{14.6}{mm.\text{年}^{-1}}$；礁体落后的速率为 $14.6 - 4.2 =
\qty{10.4}{mm.\text{年}^{-1}}$，历时 $16.5 - 8.2 = \qty{8.3}{ka}$，其顶部终止于 $10.4\times8.3 =
\qty{86}{m}$ 处。(d) $(12 - 4.2)\times8.3 = \qty{65}{m}$。两种情形下，礁体均仍位于最深的
生长极限之上，但所处水体暗得多：能否追上，取决于其生长随
深度减缓的方式。
\end{solution}

\begin{solution}{exo:b2:oceanic-islands:10}
(a) $S(t_1) = C(\sqrt{t_1 + \Delta} - \sqrt{t_1})$ 的导数为 $\frac C2(1/\sqrt{t_1 +
\Delta} - 1/\sqrt{t_1}) < 0$，故其随 $t_1$ 减小。(b) $350\sqrt{20} = \qty{1565}{m}$；$350\times
(\sqrt{70} - \sqrt{50}) = 350\times1.2955 = \qty{453}{m}$。(c) 否：$\sqrt t$ 在 $0$ 处有限，故
任一区间内的下沉量均为有限值，如 (b) 中的 \qty{1565}{m}。仅在瞬时 $t = 0$ 处的
速率为无穷大，此处该律不再用以细致描述海底。
\end{solution}

\begin{solution}{exo:b2:oceanic-islands:11}
(a) $30\times0.1 = \qty{3}{m}$。(b) 若其顶部在高海平面期正位于海平面，
则约为 \qty{120}{m}：\omterm{def:b2:oceanic-islands:reefs}{环礁}成为石灰岩岛，雨水溶解珊瑚的文石，岩石重结晶为方解石，并形成洞穴。
(c) $120/3 = 40$ 个循环，共 \qty{4}{Ma}。(d) 淋滤
带即这些暴露时期，被海水回归后再次生长的礁体封存于其下。
经过多个循环，海平面回到相近的高海平面期，故 $\Delta\eta$ 很小，礁顶
在每个高海平面期都重新生长至海平面，盖层厚度由下沉决定，正如
命题所述。
\end{solution}

\begin{solution}{exo:b2:oceanic-islands:12}
(a) $3200\times(e^{-50/62.8} - e^{-120/62.8}) = 3200\times(0.4510 - 0.1480) = \qty{970}{m}$。
(b) $150 + 970 = \qty{1120}{m}$。(c) $700\times1.8288 = \qty{1280}{m}$，比估计值深 \qty{160}{m}。
(d) 该\omterm{def:b2:oceanic-islands:seamount}{平顶海山}或许更早在较年轻的板块上淹没，并下沉得更多；邻近的
火山可能对其加载；其顶部可能被浪蚀削低，或被沉积物覆盖；其下的
板块年龄可能不同，淹没时的海平面亦非
今日的海平面。
\end{solution}

\begin{solution}{pb:b2:oceanic-islands:1}
\textbf{1.} $2500 + 350\sqrt{10} = 2500 + 350\times3.1623 = \qty{3607}{m}$.
\textbf{2.} $3607 + 1000 = \qty{4607}{m}$.
\textbf{3.} $V = \pi r^2h/3 = \pi\times50^2\times4.607/3 = \qty{12061}{km^3}$.
\textbf{4.} $0.85\times12\,061 = \qty{10252}{km^3}$，在 \qty{1}{Ma} 内：每年 \qty{0.0103}{km^3}，
约为今日基拉韦厄火山的十分之一。
\textbf{5.} $95\,600/12\,061 = 7.9$：约为八。
\textbf{6.} 拉斑质熔岩来自\omterm{def:b2:oceanic-islands:stages}{盾状阶段}；后\omterm{def:b2:oceanic-islands:stages}{盾状阶段}的碱性熔岩
覆于其上。
\textbf{7.} $\qty{1}{km}/\qty{1}{Ma} = \qty{1}{m}$ 每千年，比夏威夷岛
外海的速率慢 2.6 倍。
\textbf{8.} 是的：$4.2 - 1.0 = \qty{3.2}{mm.\text{年}^{-1}}$，尚有余量。
\textbf{9.} $12 + 1 - 4.2 = \qty{8.8}{mm}$ 每年。
\textbf{10.} $150/8.8 = \qty{17}{ka}$。侧翼上一处被淹没的礁台，
即夏威夷岛外海所见的那类台阶。
\textbf{11.} $350\times(\sqrt{22} - \sqrt{10}) = 350\times(4.6904 - 3.1623) = \qty{535}{m}$.
\textbf{12.} $1000 - 535 = \qty{465}{m}$，速率为 $465/12 = \qty{38.8}{m}$ 每百万年。
\textbf{13.} 起初为 $350/(2\sqrt{10}) = \qty{55.3}{m.Ma^{-1}}$；末了为 $350/(2\sqrt{22}) =
\qty{37.3}{m.Ma^{-1}}$。
\textbf{14.} $2500 + 350\times4.6904 = \qty{4142}{m}$.
\textbf{15.} 夏威夷的大岛高出约四倍，故海平面以上需移除的
岩石多得多；其\omterm{def:b2:oceanic-islands:stages}{盾状阶段}亦持续更久。
\textbf{16.} 海浪蚀退海岸、削出海蚀平台；滑坡与\omterm{def:b2:oceanic-islands:collapse}{山体
侧翼垮塌}。
\textbf{17.} \omterm{def:b2:oceanic-islands:reefs}{堡礁}，礁与火山海岸之间有泻湖。
\textbf{18.} $350\times(\sqrt{70} - \sqrt{22}) = 350\times(8.3666 - 4.6904) = \qty{1287}{m}$.
\textbf{19.} $1287/48 = \qty{26.8}{m}$ 每百万年，\qty{0.027}{mm} 每年：$4.2/0.0268 =
157$，即比礁体所能生长的速率慢。
\textbf{20.} 比帕里岛下玄武岩之上的石灰岩厚 \qty{20}{m}
（\qty{1267}{m}），比埃卢格拉布岛下的薄 \qty{118}{m}（\qty{1405}{m}）：二者大小相当。
\textbf{21.} $2500 + 350\times8.3666 = \qty{5428}{m}$.
\textbf{22.} $3200\times(e^{-70/62.8} - e^{-120/62.8}) = 3200\times(0.3280 - 0.1480) =
\qty{576}{m}$.
\textbf{23.} $150 + 576 = \qty{726}{m}$.
\textbf{24.} 佩特雷\omterm{def:b2:oceanic-islands:reefs}{环礁}的礁盖厚 \textbf{\qty{1287}{m}}，约 \qty{1.3}{km}
的礁灰岩，与埃尼威托克\omterm{def:b2:oceanic-islands:reefs}{环礁}之下的情况相同。
\end{solution}
